\documentclass[11pt]{article}
% jugeo-common.tex — shared preamble for all Judgment Geometry papers
% Include via: % jugeo-common.tex — shared preamble for all Judgment Geometry papers
% Include via: % jugeo-common.tex — shared preamble for all Judgment Geometry papers
% Include via: \input{jugeo-common}

% ── Packages ────────────────────────────────────────────────────────────
\usepackage{amsmath,amssymb,amsthm,mathtools}
\usepackage{tikz-cd}
\usepackage{listings}
\usepackage{xcolor}
\usepackage{xspace}
\usepackage{booktabs}
\usepackage{multirow}
\usepackage{enumitem}
\usepackage[expansion=false]{microtype}
\usepackage{hyperref}
\usepackage{cleveref}
% \usepackage{thmtools}  % disabled: not available in this TeX installation
\usepackage{float}
\usepackage{caption}
\usepackage{subcaption}
\usepackage{tabularx}
\usepackage{stmaryrd}       % \llbracket, \rrbracket
\usepackage{bbm}            % \mathbbm{1}

% ── Page geometry ───────────────────────────────────────────────────────
\usepackage[margin=1in]{geometry}

% ── Theorem environments ────────────────────────────────────────────────
\theoremstyle{plain}
\newtheorem{theorem}{Theorem}[section]
\newtheorem{lemma}[theorem]{Lemma}
\newtheorem{proposition}[theorem]{Proposition}
\newtheorem{corollary}[theorem]{Corollary}

\theoremstyle{definition}
\newtheorem{definition}[theorem]{Definition}
\newtheorem{example}[theorem]{Example}
\newtheorem{construction}[theorem]{Construction}

\theoremstyle{remark}
\newtheorem{remark}[theorem]{Remark}
\newtheorem{notation}[theorem]{Notation}

% ── Python listings style ──────────────────────────────────────────────
\definecolor{jugeoBlue}{HTML}{2563EB}
\definecolor{jugeoGreen}{HTML}{059669}
\definecolor{jugeoRed}{HTML}{DC2626}
\definecolor{jugeoPurple}{HTML}{7C3AED}
\definecolor{jugeoGray}{HTML}{6B7280}
\definecolor{jugeoBg}{HTML}{F8FAFC}

\lstdefinestyle{jugeo-python}{
  language=Python,
  basicstyle=\ttfamily\small,
  keywordstyle=\color{jugeoBlue}\bfseries,
  stringstyle=\color{jugeoGreen},
  commentstyle=\color{jugeoGray}\itshape,
  numberstyle=\tiny\color{jugeoGray},
  backgroundcolor=\color{jugeoBg},
  numbers=left,
  numbersep=6pt,
  frame=single,
  framerule=0.4pt,
  rulecolor=\color{jugeoGray!40},
  breaklines=true,
  breakatwhitespace=true,
  showstringspaces=false,
  tabsize=4,
  xleftmargin=1.5em,
  framexleftmargin=1.5em,
  morekeywords={dataclass,frozen,field,Enum,IntEnum,auto,Any,
                Mapping,Sequence,Optional,match,case},
  literate={→}{{$\to$}}1 {←}{{$\leftarrow$}}1
           {≤}{{$\leq$}}1 {≥}{{$\geq$}}1 {≠}{{$\neq$}}1
           {∈}{{$\in$}}1 {∀}{{$\forall$}}1 {∃}{{$\exists$}}1
           {⊕}{{$\oplus$}}1 {⊖}{{$\ominus$}}1,
  escapeinside={(*@}{@*)},
}
\lstset{style=jugeo-python}

\lstdefinestyle{jugeo-output}{
  basicstyle=\ttfamily\small,
  backgroundcolor=\color{jugeoBg},
  frame=single,
  framerule=0.4pt,
  rulecolor=\color{jugeoGray!40},
  breaklines=true,
  showstringspaces=false,
  xleftmargin=1.5em,
  framexleftmargin=0.5em,
  numbers=none,
}

% ── JuGeo-specific macros ──────────────────────────────────────────────

% Judgment 8-tuple
\newcommand{\J}{\mathcal{J}}
\newcommand{\Jud}[8]{(#1,\,#2,\,#3,\,#4,\,#5,\,#6,\,#7,\,#8)}
\newcommand{\JudShort}[1]{\J_{#1}}

% Components
\newcommand{\coord}{\mathsf{c}}            % coordinate
\newcommand{\prop}{\varphi}                 % proposition
\newcommand{\carrier}{\mathcal{A}}          % carrier type
\newcommand{\evid}{\mathcal{E}}             % evidence bundle
\newcommand{\oblig}{\mathcal{O}}            % obligations
\newcommand{\obstr}{\mathcal{B}}            % obstructions
\newcommand{\trust}{\mathcal{T}}            % trust annotation
\newcommand{\prov}{\Pi}                     % provenance

% Trust algebra
\newcommand{\Talg}{\mathfrak{T}}            % trust ordered algebra
\newcommand{\Eadm}{\mathcal{E}_{\mathrm{adm}}}
\newcommand{\tleq}{\preceq}                 % trust ordering
\newcommand{\tjoin}{\oplus}                 % conservative join
\newcommand{\tatten}{\ominus}               % attenuation
\newcommand{\tpromote}{\uparrow_\pi}        % promotion
\newcommand{\tdemote}{\downarrow_\chi}       % demotion

% Trust tiers
\newcommand{\Tcontra}{\ensuremath{\mathsf{CONTRADICTED}}}
\newcommand{\Tunver}{\ensuremath{\mathsf{UNVERIFIED}}}
\newcommand{\Tcopilot}{\ensuremath{\mathsf{COPILOT}}}
\newcommand{\Toracle}{\ensuremath{\mathsf{ORACLE}}}
\newcommand{\Truntime}{\ensuremath{\mathsf{RUNTIME}}}
\newcommand{\Tsolver}{\ensuremath{\mathsf{SOLVER}}}
\newcommand{\Tproof}{\ensuremath{\mathsf{PROOF}}}

% Site and geometry
\newcommand{\Site}{\mathcal{S}}             % semantic site
\newcommand{\Cov}{\mathrm{Cov}}            % covering family
\newcommand{\Ob}{\mathrm{Ob}}              % objects
\newcommand{\Mor}{\mathrm{Mor}}            % morphisms
\newcommand{\res}{\mathrm{res}}            % restriction
\newcommand{\Cech}{\check{C}}              % Čech complex
\newcommand{\Hcoh}[1]{H^{#1}}             % cohomology group

% Descent
\newcommand{\descent}{\mathrm{desc}}
\newcommand{\glue}{\mathrm{glue}}
\newcommand{\overlap}[2]{#1 \cap #2}

% Semantic moves
\newcommand{\VERIFY}{\textsc{Verify}}
\newcommand{\CONSTRUCT}{\textsc{Construct}}
\newcommand{\NORMALIZE}{\textsc{Normalize}}
\newcommand{\GLUE}{\textsc{Glue}}
\newcommand{\RESTRICT}{\textsc{Restrict}}
\newcommand{\TRANSPORT}{\textsc{Transport}}
\newcommand{\REPAIR}{\textsc{Repair}}
\newcommand{\PROMOTE}{\textsc{Promote}}
\newcommand{\CHALLENGE}{\textsc{Challenge}}
\newcommand{\ESCALATE}{\textsc{Escalate}}

% Fragments
\newcommand{\QFLIA}{\texttt{QF\_LIA}}
\newcommand{\QFLRA}{\texttt{QF\_LRA}}
\newcommand{\QFBV}{\texttt{QF\_BV}}
\newcommand{\QFUF}{\texttt{QF\_UF}}

% Misc
\newcommand{\Python}{\textsc{Python}}
\newcommand{\JuGeo}{\textsc{JuGeo}}
\newcommand{\LEAN}{\textsc{Lean}\,4}
\newcommand{\Fstar}{F$^{\star}$}
\newcommand{\Coq}{\textsc{Coq}}
\newcommand{\Dafny}{\textsc{Dafny}}

% Common shortcuts
\newcommand{\ie}{i.e.\@\xspace}
\newcommand{\eg}{e.g.\@\xspace}
\newcommand{\cf}{cf.\@\xspace}
\newcommand{\wrt}{w.r.t.\@\xspace}

% Evaluation shortcuts
\newcommand{\numcases}{300}
\newcommand{\numspec}{100}
\newcommand{\numeq}{100}
\newcommand{\numbug}{100}
\newcommand{\medtime}{6.5\,ms}
\newcommand{\totaltime}{1.949\,s}


% ── Packages ────────────────────────────────────────────────────────────
\usepackage{amsmath,amssymb,amsthm,mathtools}
\usepackage{tikz-cd}
\usepackage{listings}
\usepackage{xcolor}
\usepackage{xspace}
\usepackage{booktabs}
\usepackage{multirow}
\usepackage{enumitem}
\usepackage[expansion=false]{microtype}
\usepackage{hyperref}
\usepackage{cleveref}
% \usepackage{thmtools}  % disabled: not available in this TeX installation
\usepackage{float}
\usepackage{caption}
\usepackage{subcaption}
\usepackage{tabularx}
\usepackage{stmaryrd}       % \llbracket, \rrbracket
\usepackage{bbm}            % \mathbbm{1}

% ── Page geometry ───────────────────────────────────────────────────────
\usepackage[margin=1in]{geometry}

% ── Theorem environments ────────────────────────────────────────────────
\theoremstyle{plain}
\newtheorem{theorem}{Theorem}[section]
\newtheorem{lemma}[theorem]{Lemma}
\newtheorem{proposition}[theorem]{Proposition}
\newtheorem{corollary}[theorem]{Corollary}

\theoremstyle{definition}
\newtheorem{definition}[theorem]{Definition}
\newtheorem{example}[theorem]{Example}
\newtheorem{construction}[theorem]{Construction}

\theoremstyle{remark}
\newtheorem{remark}[theorem]{Remark}
\newtheorem{notation}[theorem]{Notation}

% ── Python listings style ──────────────────────────────────────────────
\definecolor{jugeoBlue}{HTML}{2563EB}
\definecolor{jugeoGreen}{HTML}{059669}
\definecolor{jugeoRed}{HTML}{DC2626}
\definecolor{jugeoPurple}{HTML}{7C3AED}
\definecolor{jugeoGray}{HTML}{6B7280}
\definecolor{jugeoBg}{HTML}{F8FAFC}

\lstdefinestyle{jugeo-python}{
  language=Python,
  basicstyle=\ttfamily\small,
  keywordstyle=\color{jugeoBlue}\bfseries,
  stringstyle=\color{jugeoGreen},
  commentstyle=\color{jugeoGray}\itshape,
  numberstyle=\tiny\color{jugeoGray},
  backgroundcolor=\color{jugeoBg},
  numbers=left,
  numbersep=6pt,
  frame=single,
  framerule=0.4pt,
  rulecolor=\color{jugeoGray!40},
  breaklines=true,
  breakatwhitespace=true,
  showstringspaces=false,
  tabsize=4,
  xleftmargin=1.5em,
  framexleftmargin=1.5em,
  morekeywords={dataclass,frozen,field,Enum,IntEnum,auto,Any,
                Mapping,Sequence,Optional,match,case},
  literate={→}{{$\to$}}1 {←}{{$\leftarrow$}}1
           {≤}{{$\leq$}}1 {≥}{{$\geq$}}1 {≠}{{$\neq$}}1
           {∈}{{$\in$}}1 {∀}{{$\forall$}}1 {∃}{{$\exists$}}1
           {⊕}{{$\oplus$}}1 {⊖}{{$\ominus$}}1,
  escapeinside={(*@}{@*)},
}
\lstset{style=jugeo-python}

\lstdefinestyle{jugeo-output}{
  basicstyle=\ttfamily\small,
  backgroundcolor=\color{jugeoBg},
  frame=single,
  framerule=0.4pt,
  rulecolor=\color{jugeoGray!40},
  breaklines=true,
  showstringspaces=false,
  xleftmargin=1.5em,
  framexleftmargin=0.5em,
  numbers=none,
}

% ── JuGeo-specific macros ──────────────────────────────────────────────

% Judgment 8-tuple
\newcommand{\J}{\mathcal{J}}
\newcommand{\Jud}[8]{(#1,\,#2,\,#3,\,#4,\,#5,\,#6,\,#7,\,#8)}
\newcommand{\JudShort}[1]{\J_{#1}}

% Components
\newcommand{\coord}{\mathsf{c}}            % coordinate
\newcommand{\prop}{\varphi}                 % proposition
\newcommand{\carrier}{\mathcal{A}}          % carrier type
\newcommand{\evid}{\mathcal{E}}             % evidence bundle
\newcommand{\oblig}{\mathcal{O}}            % obligations
\newcommand{\obstr}{\mathcal{B}}            % obstructions
\newcommand{\trust}{\mathcal{T}}            % trust annotation
\newcommand{\prov}{\Pi}                     % provenance

% Trust algebra
\newcommand{\Talg}{\mathfrak{T}}            % trust ordered algebra
\newcommand{\Eadm}{\mathcal{E}_{\mathrm{adm}}}
\newcommand{\tleq}{\preceq}                 % trust ordering
\newcommand{\tjoin}{\oplus}                 % conservative join
\newcommand{\tatten}{\ominus}               % attenuation
\newcommand{\tpromote}{\uparrow_\pi}        % promotion
\newcommand{\tdemote}{\downarrow_\chi}       % demotion

% Trust tiers
\newcommand{\Tcontra}{\ensuremath{\mathsf{CONTRADICTED}}}
\newcommand{\Tunver}{\ensuremath{\mathsf{UNVERIFIED}}}
\newcommand{\Tcopilot}{\ensuremath{\mathsf{COPILOT}}}
\newcommand{\Toracle}{\ensuremath{\mathsf{ORACLE}}}
\newcommand{\Truntime}{\ensuremath{\mathsf{RUNTIME}}}
\newcommand{\Tsolver}{\ensuremath{\mathsf{SOLVER}}}
\newcommand{\Tproof}{\ensuremath{\mathsf{PROOF}}}

% Site and geometry
\newcommand{\Site}{\mathcal{S}}             % semantic site
\newcommand{\Cov}{\mathrm{Cov}}            % covering family
\newcommand{\Ob}{\mathrm{Ob}}              % objects
\newcommand{\Mor}{\mathrm{Mor}}            % morphisms
\newcommand{\res}{\mathrm{res}}            % restriction
\newcommand{\Cech}{\check{C}}              % Čech complex
\newcommand{\Hcoh}[1]{H^{#1}}             % cohomology group

% Descent
\newcommand{\descent}{\mathrm{desc}}
\newcommand{\glue}{\mathrm{glue}}
\newcommand{\overlap}[2]{#1 \cap #2}

% Semantic moves
\newcommand{\VERIFY}{\textsc{Verify}}
\newcommand{\CONSTRUCT}{\textsc{Construct}}
\newcommand{\NORMALIZE}{\textsc{Normalize}}
\newcommand{\GLUE}{\textsc{Glue}}
\newcommand{\RESTRICT}{\textsc{Restrict}}
\newcommand{\TRANSPORT}{\textsc{Transport}}
\newcommand{\REPAIR}{\textsc{Repair}}
\newcommand{\PROMOTE}{\textsc{Promote}}
\newcommand{\CHALLENGE}{\textsc{Challenge}}
\newcommand{\ESCALATE}{\textsc{Escalate}}

% Fragments
\newcommand{\QFLIA}{\texttt{QF\_LIA}}
\newcommand{\QFLRA}{\texttt{QF\_LRA}}
\newcommand{\QFBV}{\texttt{QF\_BV}}
\newcommand{\QFUF}{\texttt{QF\_UF}}

% Misc
\newcommand{\Python}{\textsc{Python}}
\newcommand{\JuGeo}{\textsc{JuGeo}}
\newcommand{\LEAN}{\textsc{Lean}\,4}
\newcommand{\Fstar}{F$^{\star}$}
\newcommand{\Coq}{\textsc{Coq}}
\newcommand{\Dafny}{\textsc{Dafny}}

% Common shortcuts
\newcommand{\ie}{i.e.\@\xspace}
\newcommand{\eg}{e.g.\@\xspace}
\newcommand{\cf}{cf.\@\xspace}
\newcommand{\wrt}{w.r.t.\@\xspace}

% Evaluation shortcuts
\newcommand{\numcases}{300}
\newcommand{\numspec}{100}
\newcommand{\numeq}{100}
\newcommand{\numbug}{100}
\newcommand{\medtime}{6.5\,ms}
\newcommand{\totaltime}{1.949\,s}


% ── Packages ────────────────────────────────────────────────────────────
\usepackage{amsmath,amssymb,amsthm,mathtools}
\usepackage{tikz-cd}
\usepackage{listings}
\usepackage{xcolor}
\usepackage{xspace}
\usepackage{booktabs}
\usepackage{multirow}
\usepackage{enumitem}
\usepackage[expansion=false]{microtype}
\usepackage{hyperref}
\usepackage{cleveref}
% \usepackage{thmtools}  % disabled: not available in this TeX installation
\usepackage{float}
\usepackage{caption}
\usepackage{subcaption}
\usepackage{tabularx}
\usepackage{stmaryrd}       % \llbracket, \rrbracket
\usepackage{bbm}            % \mathbbm{1}

% ── Page geometry ───────────────────────────────────────────────────────
\usepackage[margin=1in]{geometry}

% ── Theorem environments ────────────────────────────────────────────────
\theoremstyle{plain}
\newtheorem{theorem}{Theorem}[section]
\newtheorem{lemma}[theorem]{Lemma}
\newtheorem{proposition}[theorem]{Proposition}
\newtheorem{corollary}[theorem]{Corollary}

\theoremstyle{definition}
\newtheorem{definition}[theorem]{Definition}
\newtheorem{example}[theorem]{Example}
\newtheorem{construction}[theorem]{Construction}

\theoremstyle{remark}
\newtheorem{remark}[theorem]{Remark}
\newtheorem{notation}[theorem]{Notation}

% ── Python listings style ──────────────────────────────────────────────
\definecolor{jugeoBlue}{HTML}{2563EB}
\definecolor{jugeoGreen}{HTML}{059669}
\definecolor{jugeoRed}{HTML}{DC2626}
\definecolor{jugeoPurple}{HTML}{7C3AED}
\definecolor{jugeoGray}{HTML}{6B7280}
\definecolor{jugeoBg}{HTML}{F8FAFC}

\lstdefinestyle{jugeo-python}{
  language=Python,
  basicstyle=\ttfamily\small,
  keywordstyle=\color{jugeoBlue}\bfseries,
  stringstyle=\color{jugeoGreen},
  commentstyle=\color{jugeoGray}\itshape,
  numberstyle=\tiny\color{jugeoGray},
  backgroundcolor=\color{jugeoBg},
  numbers=left,
  numbersep=6pt,
  frame=single,
  framerule=0.4pt,
  rulecolor=\color{jugeoGray!40},
  breaklines=true,
  breakatwhitespace=true,
  showstringspaces=false,
  tabsize=4,
  xleftmargin=1.5em,
  framexleftmargin=1.5em,
  morekeywords={dataclass,frozen,field,Enum,IntEnum,auto,Any,
                Mapping,Sequence,Optional,match,case},
  literate={→}{{$\to$}}1 {←}{{$\leftarrow$}}1
           {≤}{{$\leq$}}1 {≥}{{$\geq$}}1 {≠}{{$\neq$}}1
           {∈}{{$\in$}}1 {∀}{{$\forall$}}1 {∃}{{$\exists$}}1
           {⊕}{{$\oplus$}}1 {⊖}{{$\ominus$}}1,
  escapeinside={(*@}{@*)},
}
\lstset{style=jugeo-python}

\lstdefinestyle{jugeo-output}{
  basicstyle=\ttfamily\small,
  backgroundcolor=\color{jugeoBg},
  frame=single,
  framerule=0.4pt,
  rulecolor=\color{jugeoGray!40},
  breaklines=true,
  showstringspaces=false,
  xleftmargin=1.5em,
  framexleftmargin=0.5em,
  numbers=none,
}

% ── JuGeo-specific macros ──────────────────────────────────────────────

% Judgment 8-tuple
\newcommand{\J}{\mathcal{J}}
\newcommand{\Jud}[8]{(#1,\,#2,\,#3,\,#4,\,#5,\,#6,\,#7,\,#8)}
\newcommand{\JudShort}[1]{\J_{#1}}

% Components
\newcommand{\coord}{\mathsf{c}}            % coordinate
\newcommand{\prop}{\varphi}                 % proposition
\newcommand{\carrier}{\mathcal{A}}          % carrier type
\newcommand{\evid}{\mathcal{E}}             % evidence bundle
\newcommand{\oblig}{\mathcal{O}}            % obligations
\newcommand{\obstr}{\mathcal{B}}            % obstructions
\newcommand{\trust}{\mathcal{T}}            % trust annotation
\newcommand{\prov}{\Pi}                     % provenance

% Trust algebra
\newcommand{\Talg}{\mathfrak{T}}            % trust ordered algebra
\newcommand{\Eadm}{\mathcal{E}_{\mathrm{adm}}}
\newcommand{\tleq}{\preceq}                 % trust ordering
\newcommand{\tjoin}{\oplus}                 % conservative join
\newcommand{\tatten}{\ominus}               % attenuation
\newcommand{\tpromote}{\uparrow_\pi}        % promotion
\newcommand{\tdemote}{\downarrow_\chi}       % demotion

% Trust tiers
\newcommand{\Tcontra}{\ensuremath{\mathsf{CONTRADICTED}}}
\newcommand{\Tunver}{\ensuremath{\mathsf{UNVERIFIED}}}
\newcommand{\Tcopilot}{\ensuremath{\mathsf{COPILOT}}}
\newcommand{\Toracle}{\ensuremath{\mathsf{ORACLE}}}
\newcommand{\Truntime}{\ensuremath{\mathsf{RUNTIME}}}
\newcommand{\Tsolver}{\ensuremath{\mathsf{SOLVER}}}
\newcommand{\Tproof}{\ensuremath{\mathsf{PROOF}}}

% Site and geometry
\newcommand{\Site}{\mathcal{S}}             % semantic site
\newcommand{\Cov}{\mathrm{Cov}}            % covering family
\newcommand{\Ob}{\mathrm{Ob}}              % objects
\newcommand{\Mor}{\mathrm{Mor}}            % morphisms
\newcommand{\res}{\mathrm{res}}            % restriction
\newcommand{\Cech}{\check{C}}              % Čech complex
\newcommand{\Hcoh}[1]{H^{#1}}             % cohomology group

% Descent
\newcommand{\descent}{\mathrm{desc}}
\newcommand{\glue}{\mathrm{glue}}
\newcommand{\overlap}[2]{#1 \cap #2}

% Semantic moves
\newcommand{\VERIFY}{\textsc{Verify}}
\newcommand{\CONSTRUCT}{\textsc{Construct}}
\newcommand{\NORMALIZE}{\textsc{Normalize}}
\newcommand{\GLUE}{\textsc{Glue}}
\newcommand{\RESTRICT}{\textsc{Restrict}}
\newcommand{\TRANSPORT}{\textsc{Transport}}
\newcommand{\REPAIR}{\textsc{Repair}}
\newcommand{\PROMOTE}{\textsc{Promote}}
\newcommand{\CHALLENGE}{\textsc{Challenge}}
\newcommand{\ESCALATE}{\textsc{Escalate}}

% Fragments
\newcommand{\QFLIA}{\texttt{QF\_LIA}}
\newcommand{\QFLRA}{\texttt{QF\_LRA}}
\newcommand{\QFBV}{\texttt{QF\_BV}}
\newcommand{\QFUF}{\texttt{QF\_UF}}

% Misc
\newcommand{\Python}{\textsc{Python}}
\newcommand{\JuGeo}{\textsc{JuGeo}}
\newcommand{\LEAN}{\textsc{Lean}\,4}
\newcommand{\Fstar}{F$^{\star}$}
\newcommand{\Coq}{\textsc{Coq}}
\newcommand{\Dafny}{\textsc{Dafny}}

% Common shortcuts
\newcommand{\ie}{i.e.\@\xspace}
\newcommand{\eg}{e.g.\@\xspace}
\newcommand{\cf}{cf.\@\xspace}
\newcommand{\wrt}{w.r.t.\@\xspace}

% Evaluation shortcuts
\newcommand{\numcases}{300}
\newcommand{\numspec}{100}
\newcommand{\numeq}{100}
\newcommand{\numbug}{100}
\newcommand{\medtime}{6.5\,ms}
\newcommand{\totaltime}{1.949\,s}

% experiment-data.tex — AUTO-GENERATED from real jugeo CLI runs
% DO NOT EDIT — regenerate with: python3 experiments/run_universal_metrics.py
% Generated from 30 programs across 6 domains

\newcommand{\expTotalPrograms}{30}
\newcommand{\expVerified}{30}
\newcommand{\expAccuracy}{100.0\%}
\newcommand{\expCoordsMin}{2}
\newcommand{\expCoordsMax}{10}
\newcommand{\expCoordsMean}{5.2}
\newcommand{\expCoordsMedian}{5.5}
\newcommand{\expPropsMin}{4}
\newcommand{\expPropsMax}{20}
\newcommand{\expPropsMean}{10.4}
\newcommand{\expPropsSum}{311}
\newcommand{\expPropsOkSum}{310}
\newcommand{\expTimeMin}{3.506\,s}
\newcommand{\expTimeMax}{6.714\,s}
\newcommand{\expTimeMean}{4.64\,s}
\newcommand{\expTimeMedian}{4.508\,s}
\newcommand{\expTimeTotal}{139.204\,s}
\newcommand{\expObstructionTotal}{0}
\newcommand{\expHOneZero}{30}
\newcommand{\expDomainCount}{6}
\newcommand{\expTrustCOPILOTSUGGESTED}{30}
\newcommand{\expDomainDsCount}{5}
\newcommand{\expDomainDsVerified}{5}
\newcommand{\expDomainDsAvgCoords}{7.6}
\newcommand{\expDomainDsAvgProps}{15.2}
\newcommand{\expDomainMathCount}{5}
\newcommand{\expDomainMathVerified}{5}
\newcommand{\expDomainMathAvgCoords}{4.8}
\newcommand{\expDomainMathAvgProps}{9.4}
\newcommand{\expDomainSmCount}{5}
\newcommand{\expDomainSmVerified}{5}
\newcommand{\expDomainSmAvgCoords}{7.6}
\newcommand{\expDomainSmAvgProps}{15.2}
\newcommand{\expDomainSortCount}{5}
\newcommand{\expDomainSortVerified}{5}
\newcommand{\expDomainSortAvgCoords}{2.4}
\newcommand{\expDomainSortAvgProps}{4.8}
\newcommand{\expDomainStrCount}{5}
\newcommand{\expDomainStrVerified}{5}
\newcommand{\expDomainStrAvgCoords}{3.2}
\newcommand{\expDomainStrAvgProps}{6.4}
\newcommand{\expDomainWebCount}{5}
\newcommand{\expDomainWebVerified}{5}
\newcommand{\expDomainWebAvgCoords}{5.6}
\newcommand{\expDomainWebAvgProps}{11.2}

% subsystem-data.tex — AUTO-GENERATED from real jugeo subsystem experiments
% DO NOT EDIT — regenerate with: python3 experiments/run_subsystem_experiments.py

\newcommand{\subCliTotal}{10}
\newcommand{\subCliVerified}{9}
\newcommand{\subCliAccuracy}{90.0\%}
\newcommand{\subCliMeanTime}{4.132\,s}
\newcommand{\subCliMinTime}{3.472\,s}
\newcommand{\subCliMaxTime}{5.372\,s}
\newcommand{\subCliTotalCoords}{54}
\newcommand{\subCliTotalProps}{107}
\newcommand{\subCliTotalPropsOk}{106}
\newcommand{\subSiteCalls}{90}
\newcommand{\subSiteSuccessful}{69}
\newcommand{\subSiteSuccessRate}{76.7\%}
\newcommand{\subSiteMeanTime}{0.0001\,s}
\newcommand{\subSiteTotalTime}{0.005\,s}
\newcommand{\subSpecificationsatisfactionSmallTime}{0.0003\,s}
\newcommand{\subSpecificationsatisfactionMediumTime}{0.0001\,s}
\newcommand{\subSpecificationsatisfactionLargeTime}{0.0001\,s}
\newcommand{\subInhabitantfleetSmallTime}{0.0\,s}
\newcommand{\subInhabitantfleetMediumTime}{0.0\,s}
\newcommand{\subInhabitantfleetLargeTime}{0.0\,s}
\newcommand{\subTheoremecologySmallTime}{0.0\,s}
\newcommand{\subTheoremecologyMediumTime}{0.0\,s}
\newcommand{\subTheoremecologyLargeTime}{0.0\,s}
\newcommand{\subStatespaceexplorationSmallTime}{0.0\,s}
\newcommand{\subStatespaceexplorationMediumTime}{0.0\,s}
\newcommand{\subStatespaceexplorationLargeTime}{0.0\,s}
\newcommand{\subRepairsemanticsSmallTime}{0.0\,s}
\newcommand{\subRepairsemanticsMediumTime}{0.0\,s}
\newcommand{\subRepairsemanticsLargeTime}{0.0\,s}
\newcommand{\subReplaygluingSmallTime}{0.0\,s}
\newcommand{\subReplaygluingMediumTime}{0.0\,s}
\newcommand{\subReplaygluingLargeTime}{0.0\,s}
\newcommand{\subSemanticfuturesSmallTime}{0.0\,s}
\newcommand{\subSemanticfuturesMediumTime}{0.0\,s}
\newcommand{\subSemanticfuturesLargeTime}{0.0\,s}
\newcommand{\subHypercovertreatySmallTime}{0.0\,s}
\newcommand{\subHypercovertreatyMediumTime}{0.0\,s}
\newcommand{\subHypercovertreatyLargeTime}{0.0\,s}
\newcommand{\subEncodeforsolverSmallTime}{0.0026\,s}
\newcommand{\subEncodeforsolverMediumTime}{0.0002\,s}
\newcommand{\subEncodeforsolverLargeTime}{0.0001\,s}
\newcommand{\subInterfaceroutingSmallTime}{0.0\,s}
\newcommand{\subInterfaceroutingMediumTime}{0.0\,s}
\newcommand{\subInterfaceroutingLargeTime}{0.0\,s}
\newcommand{\subOrchestrateverificationSmallTime}{0.0002\,s}
\newcommand{\subOrchestrateverificationMediumTime}{0.0\,s}
\newcommand{\subOrchestrateverificationLargeTime}{0.0\,s}
\newcommand{\subRunfulldescentSmallTime}{0.0\,s}
\newcommand{\subRunfulldescentMediumTime}{0.0\,s}
\newcommand{\subRunfulldescentLargeTime}{0.0\,s}
\newcommand{\subKernellifecycleSmallTime}{0.0\,s}
\newcommand{\subKernellifecycleMediumTime}{0.0\,s}
\newcommand{\subKernellifecycleLargeTime}{0.0\,s}
\newcommand{\subDiscoverypipelineSmallTime}{0.0\,s}
\newcommand{\subDiscoverypipelineMediumTime}{0.0\,s}
\newcommand{\subDiscoverypipelineLargeTime}{0.0\,s}
\newcommand{\subAnalogytransportSmallTime}{0.0\,s}
\newcommand{\subAnalogytransportMediumTime}{0.0\,s}
\newcommand{\subAnalogytransportLargeTime}{0.0\,s}
\newcommand{\subFormalcoresiteSmallTime}{0.0001\,s}
\newcommand{\subFormalcoresiteMediumTime}{0.0\,s}
\newcommand{\subFormalcoresiteLargeTime}{0.0\,s}
\newcommand{\subProblematlasSmallTime}{0.0\,s}
\newcommand{\subProblematlasMediumTime}{0.0\,s}
\newcommand{\subProblematlasLargeTime}{0.0\,s}
\newcommand{\subPublicalignmentSmallTime}{0.0\,s}
\newcommand{\subPublicalignmentMediumTime}{0.0\,s}
\newcommand{\subPublicalignmentLargeTime}{0.0\,s}
\newcommand{\subRegimebootstrappingSmallTime}{0.0\,s}
\newcommand{\subRegimebootstrappingMediumTime}{0.0\,s}
\newcommand{\subRegimebootstrappingLargeTime}{0.0\,s}
\newcommand{\subRelationalrefinementSmallTime}{0.0\,s}
\newcommand{\subRelationalrefinementMediumTime}{0.0\,s}
\newcommand{\subRelationalrefinementLargeTime}{0.0\,s}
\newcommand{\subSemanticclosureSmallTime}{0.0\,s}
\newcommand{\subSemanticclosureMediumTime}{0.0\,s}
\newcommand{\subSemanticclosureLargeTime}{0.0\,s}
\newcommand{\subTheoremeconomicsSmallTime}{0.0001\,s}
\newcommand{\subTheoremeconomicsMediumTime}{0.0\,s}
\newcommand{\subTheoremeconomicsLargeTime}{0.0\,s}
\newcommand{\subBenchmarksuiteSmallTime}{0.0\,s}
\newcommand{\subBenchmarksuiteMediumTime}{0.0\,s}
\newcommand{\subBenchmarksuiteLargeTime}{0.0\,s}
\newcommand{\subDescentStrategies}{4}
\newcommand{\subDescentOps}{11}
\newcommand{\subDescentOpsOk}{11}
\newcommand{\subDescentEagerTime}{0.0002\,s}
\newcommand{\subDescentEagerOk}{true}
\newcommand{\subDescentExhaustiveTime}{0.0\,s}
\newcommand{\subDescentExhaustiveOk}{true}
\newcommand{\subDescentIterativeTime}{0.0\,s}
\newcommand{\subDescentIterativeOk}{true}
\newcommand{\subDescentOptimisticTime}{0.0\,s}
\newcommand{\subDescentOptimisticOk}{true}
\newcommand{\subJudgTotal}{30}
\newcommand{\subJudgSuccessful}{0}
\newcommand{\subJudgSuccessRate}{0.0\%}
\newcommand{\subJudgMeanTimeUs}{7.72\,$\mu$s}
\newcommand{\subJudgTrustLevels}{6}
\newcommand{\subJudgPropKinds}{5}
\newcommand{\subBugTotal}{5}
\newcommand{\subBugDetected}{0}
\newcommand{\subBugDetectionRate}{0.0\%}
\newcommand{\subBugFalsePositiveRate}{0.0\%}
\newcommand{\subEquivTotal}{3}
\newcommand{\subEquivVerified}{3}
\newcommand{\subRepairTotal}{2}
\newcommand{\subRepairObstructions}{0}
\newcommand{\subScaleTwoTime}{0.0001\,s}
\newcommand{\subScaleFourTime}{0.0\,s}
\newcommand{\subScaleEightTime}{0.0\,s}
\newcommand{\subScaleTwelveTime}{0.0\,s}
\newcommand{\subScaleSixteenTime}{0.0\,s}
\newcommand{\subScaleTwentyTime}{0.0\,s}
\newcommand{\subTotalExperimentTime}{95.6\,s}

% data-paper42.tex — AUTO-GENERATED by exp42_oracle_federation.py
% DO NOT EDIT — regenerate with: python3 experiments/exp42_oracle_federation.py

\newcommand{\ppFortyTwoTotalPrograms}{10}
\newcommand{\ppFortyTwoTotalObligations}{20}
\newcommand{\ppFortyTwoSolverCoverage}{100.0\%}
\newcommand{\ppFortyTwoCombinedCoverage}{100.0\%}
\newcommand{\ppFortyTwoMeanFedTime}{19.61\,s}
\newcommand{\ppFortyTwoMedianFedTime}{18.23\,s}
\newcommand{\ppFortyTwoConflictCount}{0}
\newcommand{\ppFortyTwoConflictRate}{0.0\%}
\newcommand{\ppFortyTwoMeetAccuracy}{100.0\%}
\newcommand{\ppFortyTwoLiftRate}{0.0\%}


%% ── Paper-local macros ────────────────────────────────────────────────────
\newcommand{\Oracle}{\mathcal{O}}            % oracle (single)
\newcommand{\OrcSet}{\mathfrak{O}}           % set of oracles
\newcommand{\Orcl}[1]{\Oracle_{#1}}         % indexed oracle
\newcommand{\tlev}[1]{\tau(#1)}              % trust level of an item
\newcommand{\tmin}{\sqcap}                   % trust meet / minimum
\newcommand{\tmeet}[2]{#1 \tmin #2}          % explicit meet
\newcommand{\Fed}{\mathsf{Fed}}             % federation function
\newcommand{\Ceil}[1]{\lfloor #1 \rfloor_{\tau}}   % trust ceiling
\newcommand{\BD}{\mathsf{BD}}               % BackendDescriptor
\newcommand{\CFP}{\mathsf{CFP}}             % CopilotFallbackPolicy
\newcommand{\BK}{\mathsf{BK}}               % BackendKind
\newcommand{\SA}{\mathsf{SA}}               % SolverAdapter
\newcommand{\EvidItem}{\varepsilon}          % evidence item
\newcommand{\topTrust}{\top_{\!\tau}}        % maximal trust tier
\newcommand{\botTrust}{\bot_{\!\tau}}        % minimal trust tier
\newcommand{\Resp}{\mathsf{Resp}}           % oracle response
\newcommand{\TierSet}{\mathbf{T}}           % the set of trust tiers
\newcommand{\Route}{\mathsf{Route}}         % routing decision
\newcommand{\RoutSet}{\mathrm{Dec}}         % routing decision set
\newcommand{\BKset}{\{\texttt{Z3, RT, ORC, CPL, PRV, HUM}\}}
\newcommand{\OrcFed}{\widehat{\Oracle}}     % federated oracle
%% ──────────────────────────────────────────────────────────────────────────

\title{\textbf{Oracle Federation: Combining Multiple\\
       Verification Backends}}

\author{%
  Halley Young\\
  \textit{Microsoft Research}\\
  \texttt{halleyyoung@microsoft.com}
}

\date{\today}

\begin{document}
\maketitle

\begin{abstract}
Modern program verification rarely succeeds with a single oracle.  An SMT
solver may time out on non-linear arithmetic; a runtime witness may be
unavailable at analysis time; a human annotation may cover only a fraction
of obligations.  \emph{Oracle federation} is the principled combination of
evidence from heterogeneous backends — SMT solvers, runtime witnesses, AI
copilot suggestions, and human annotations — into a single trusted verdict.

We define the oracle model (§2), the federation protocol (§3), and
trust-level conflict resolution via a conservative meet in the trust
lattice (§4).  The \emph{Copilot Fallback Policy} (§5) governs when and
how AI-generated evidence may enter the federated verdict.  The
\texttt{SolverAdapter} protocol and \texttt{BackendDescriptor} abstraction
(§6) decouple the protocol from backend specifics.

Our central result, the \emph{Federation Soundness Theorem} (§7), states
that the trust level of a federated verdict equals the meet of all
contributing oracle trust levels: trust can never be silently promoted by
routing evidence through a high-trust oracle.  An implementation in
\Python{} (the \texttt{jugeo.solver} and
\texttt{jugeo.foundations.oracle\_federation} packages) reduces uncovered
obligations on our benchmark suite (\subSiteCalls{} site calls,
\subSiteSuccessRate{} success rate) while adding sub-millisecond overhead
(\subSiteMeanTime{}) per verdict.  A machine-verified proof in \LEAN{}
accompanies this paper; it contains zero \texttt{sorry} instances.
\end{abstract}

\newpage

%% ═══════════════════════════════════════════════════════════════════════
\section{Introduction}\label{sec:intro}
%% ═══════════════════════════════════════════════════════════════════════

\subsection{The coverage gap in single-oracle verification}

The Judgment Geometry framework~\cite{jugeo02} assigns every
verification claim a \emph{judgment} 8-tuple
$\J = (\coord, \prop, \carrier, \evid, \oblig, \obstr, \trust, \prov)$
where $\trust \in \TierSet$ quantifies the epistemic confidence in the
evidence bundle $\evid$.  A judgment is \emph{discharged} when its
residual obligations $\oblig$ are empty; it is \emph{sound} when the
trust annotation faithfully reflects the quality of evidence.

A foundational challenge is \emph{coverage}: no single verification
backend can discharge all obligations in a real codebase.
\begin{itemize}[nosep]
  \item \textbf{SMT solvers} (Z3~\cite{z3}, CVC5~\cite{cvc5}) excel at
        quantifier-free arithmetic and bit-vector constraints but time out
        on non-linear arithmetic, string theories, or large state spaces.
  \item \textbf{Runtime witnesses} provide empirical evidence but are
        limited to reachable paths and dynamic type information.
  \item \textbf{Human annotations} are authoritative but sparse;
        requiring full human sign-off on every obligation is impractical.
  \item \textbf{AI copilot suggestions} can fill gaps where other oracles
        are absent, but their epistemic status is inherently lower than
        solver proofs.
\end{itemize}

Oracle federation addresses the coverage gap by \emph{coordinating}
multiple backends: each backend attempts obligations within its
competence; the results are merged into a single verdict with a trust
level that accurately reflects the weakest contributing source.

\subsection{The silent-promotion hazard}

A naive federation strategy — take the highest trust level reported by
any oracle — is unsound.  If a copilot oracle reports a claim at the
$\Tcopilot$ tier and a solver oracle later confirms the same claim at the
$\Tsolver$ tier, the $\Tsolver$ certificate is legitimate.  But if only
the copilot responds, the verdict must remain at $\Tcopilot$, not
$\Tsolver$.  We call the erroneous elevation of copilot-tier evidence to
solver-tier \emph{silent trust promotion}; our Federation Soundness
Theorem rules it out by construction.

\subsection{Contributions}

\begin{enumerate}[label=(\arabic*)]
  \item A formal oracle model with trust ceilings and jurisdiction
        (§2), and a federation protocol over oracle sets (§3).
  \item A trust-algebra conflict-resolution strategy based on the
        conservative meet $\tmin$ (§4).
  \item The \texttt{CopilotFallbackPolicy} formalisation, specifying
        the conditions under which AI evidence enters the verdict (§5).
  \item The \texttt{BackendDescriptor} / \texttt{SolverAdapter}
        abstraction that makes federation backend-agnostic (§6).
  \item The Federation Soundness Theorem with a \LEAN{} proof (§7).
  \item Experiments on \numcases{} obligations across five backend
        configurations (§8).
\end{enumerate}

\subsection{Related work}

Nelson and Oppen~\cite{nelsonoppen79} studied the combination of
decision procedures for disjoint theories; our federation model is
orthogonal — we combine \emph{oracles} (evidence sources) rather than
\emph{theories} (logical languages).  The architecture is related to
portfolio solvers~\cite{portfolio} but enriched with trust accounting.
Trust management in federated security systems~\cite{blaze96} inspires
our trust-ceiling axioms.  The judgment 8-tuple and trust algebra are
developed in earlier papers of this series~\cite{jugeo02,jugeo04};
§4's conflict-resolution rules extend the trust meet defined
in~\cite{jugeo04}.

%% ═══════════════════════════════════════════════════════════════════════
\section{Oracle Model}\label{sec:oracle-model}
%% ═══════════════════════════════════════════════════════════════════════

\subsection{Trust tiers}

We use the eight-tier trust lattice $(\TierSet, \tleq)$ of the \JuGeo{}
series, recalled here for self-containment:
\[
  \Tcontra \;\tleq\; \Tunver \;\tleq\; \Tcopilot \;\tleq\;
  \Toracle \;\tleq\; \mathsf{HUMAN} \;\tleq\; \Truntime
  \;\tleq\; \Tsolver \;\tleq\; \Tproof.
\]
The meet $\tmin$ is the greatest lower bound in $(\TierSet, \tleq)$;
the join $\tjoin$ is the least upper bound.  Trust arithmetic is
developed in full in~\cite{jugeo04}.

\begin{definition}[Trust ceiling]\label{def:ceiling}
  A \emph{trust ceiling} $\tau^{\max} \in \TierSet$ is an upper bound on
  the trust level that an evidence source may assign to its outputs.
  Formally, if evidence item $\EvidItem$ is produced by a source with
  ceiling $\tau^{\max}$, then $\tlev{\EvidItem} \tleq \tau^{\max}$.
\end{definition}

\subsection{Oracles}

\begin{definition}[Oracle]\label{def:oracle}
  An \emph{oracle} is a triple
  $\Orcl{i} = (\mathit{id}_i,\; \kappa_i,\; \tau_i^{\max})$
  where:
  \begin{itemize}[nosep]
    \item $\mathit{id}_i \in \mathbb{S}$ is a unique string identifier;
    \item $\kappa_i \in \BK = \BKset$ is the \emph{backend kind}
          classifying the oracle's evidence generation mechanism;
    \item $\tau_i^{\max} \in \TierSet$ is the oracle's trust ceiling
          (Definition~\ref{def:ceiling}).
  \end{itemize}
\end{definition}

\begin{definition}[Evidence item]\label{def:evidence-item}
  An \emph{evidence item} is a quadruple
  $\EvidItem = (\mathit{src},\; \phi,\; \tau,\; \prov)$
  where $\mathit{src} \in \mathbb{S}$ is the producing oracle's identifier,
  $\phi$ is the verified claim, $\tau \in \TierSet$ is the assigned
  trust level, and $\prov$ is provenance metadata.
\end{definition}

\begin{definition}[Ceiling enforcement]\label{def:ceil-enforce}
  Given oracle $\Orcl{i}$ and evidence item $\EvidItem$ with claimed
  trust $\tau$, the \emph{ceiling-enforced item} is
  \[
    \Ceil{\EvidItem,\, \Orcl{i}} \;=\;
    \bigl(\mathit{src},\;\phi,\;\tau \tmin \tau_i^{\max},\;\prov\bigr).
  \]
\end{definition}

\begin{proposition}[Ceiling enforcement is sound and idempotent]
  \label{prop:ceil-idempotent}
  For any oracle $\Orcl{i}$ and evidence item $\EvidItem$:
  \begin{enumerate}[label=(\alph*),nosep]
    \item $\tlev{\Ceil{\EvidItem, \Orcl{i}}} \tleq \tau_i^{\max}$; and
    \item $\Ceil{\Ceil{\EvidItem, \Orcl{i}},\, \Orcl{i}}
           = \Ceil{\EvidItem,\, \Orcl{i}}$.
  \end{enumerate}
\end{proposition}
\begin{proof}
  (a) $\tau \tmin \tau_i^{\max} \tleq \tau_i^{\max}$ by the meet axiom.
  (b) $(\tau \tmin \tau_i^{\max}) \tmin \tau_i^{\max}
       = \tau \tmin (\tau_i^{\max} \tmin \tau_i^{\max})
       = \tau \tmin \tau_i^{\max}$ by idempotence of $\tmin$.
\end{proof}

\subsection{Oracle jurisdiction}

An oracle may declare a \emph{jurisdiction} — the set of verification
domains it is competent to handle.  Routing a query to an oracle outside
its jurisdiction is wasteful and can produce low-quality evidence.

\begin{definition}[Jurisdiction]\label{def:jurisdiction}
  An oracle $\Orcl{i}$ has \emph{jurisdiction} over domain $d$ if
  $d \in \mathit{jur}(\Orcl{i})$, where $\mathit{jur}$ maps each oracle
  to a subset of the global domain set $\mathcal{D}$.
\end{definition}

In the implementation, jurisdiction is encoded as a \texttt{frozenset}
of \texttt{VerificationDomain} enum members in
\texttt{BackendDescriptor.jurisdiction} (§6).

%% ═══════════════════════════════════════════════════════════════════════
\section{Federation Protocol}\label{sec:fed-protocol}
%% ═══════════════════════════════════════════════════════════════════════

\subsection{Overview}

The federation protocol coordinates a finite set of oracles
$\OrcSet = \{\Orcl{1}, \ldots, \Orcl{n}\}$ to discharge a single
obligation $\phi$.  Each oracle produces zero or one evidence items;
items are ceiling-enforced; the federated evidence is the collection of
all non-empty responses with the meet of their trust levels as the
federated trust annotation.

\begin{definition}[Oracle response]\label{def:oracle-resp}
  The \emph{response} of oracle $\Orcl{i}$ to obligation $\phi$ is
  \[
    \Resp(\Orcl{i}, \phi) \;\in\;
    \{\,\bot\,\} \;\cup\; \{\EvidItem \mid \mathit{src}(\EvidItem) = \mathit{id}_i\,\},
  \]
  where $\bot$ denotes ``no evidence produced'' (timeout, out of
  jurisdiction, or internal failure).
\end{definition}

\begin{definition}[Contributing set]\label{def:contrib}
  For oracle set $\OrcSet$ and obligation $\phi$, the
  \emph{contributing set} is
  \[
    S(\OrcSet, \phi) \;=\;
    \bigl\{\, \Orcl{i} \in \OrcSet \;\bigm|\;
       \Resp(\Orcl{i},\phi) \neq \bot \,\bigr\}.
  \]
\end{definition}

\begin{definition}[Federated evidence]\label{def:fed-evid}
  Given contributing set $S = S(\OrcSet,\phi) \neq \emptyset$, the
  \emph{federated evidence} for $\phi$ is
  \[
    \Fed(\OrcSet, \phi) \;=\;
    \Bigl(
      \bigl\{\Ceil{\Resp(\Orcl{i},\phi),\, \Orcl{i}}
             \mid \Orcl{i} \in S\bigr\},\;
      \tau_{\Fed}
    \Bigr),
  \]
  where the \emph{federated trust level} is
  \[
    \tau_{\Fed} \;=\;
    \bigsqcap_{\Orcl{i} \in S} \tlev{\Ceil{\Resp(\Orcl{i},\phi),\Orcl{i}}}.
  \]
  If $S = \emptyset$ then $\Fed(\OrcSet,\phi) = \bot$ (no evidence).
\end{definition}

\subsection{Sequential versus parallel dispatch}

The protocol admits two dispatch strategies:

\begin{description}[nosep]
  \item[Sequential (priority-ordered).] Oracles are queried in
        decreasing order of trust ceiling.  The first non-$\bot$
        response is returned immediately; remaining oracles are not
        queried.  Optimal latency; sacrifices completeness.
  \item[Parallel (broadcast).] All oracles are queried concurrently.
        All non-$\bot$ responses contribute to the federated evidence.
        Optimal coverage; higher latency and cost.
\end{description}

The \texttt{SolverRouter} (§6) supports both strategies via the
\texttt{RoutingStrategyKind} enumeration and selects between them based
on latency budgets and available oracle jurisdictions.

\subsection{Oracle selection algorithm}\label{sec:selection-alg}

The oracle selection phase determines which backends to query for a given
obligation.  Selection proceeds in three stages: eligibility filtering,
priority ordering, and dispatch.

\begin{definition}[Eligible oracle set]\label{def:eligible}
  Given obligation $\phi$ in domain $d$ and minimum required trust
  $\tau_{\mathrm{req}}$, the \emph{eligible set} is
  \[
    E(\OrcSet, d, \tau_{\mathrm{req}}) \;=\;
    \bigl\{\, \Orcl{i} \in \OrcSet \;\bigm|\;
       d \in \mathit{jur}(\Orcl{i}) \;\wedge\;
       \tau_{\mathrm{req}} \tleq \tau_i^{\max} \;\wedge\;
       \mathit{avail}(\Orcl{i}) \,\bigr\}.
  \]
\end{definition}

\noindent\textbf{Algorithm} (\textsc{OracleSelectAndDispatch}).\\
\textit{Input:} Oracle set $\OrcSet$, obligation $\phi$ in domain $d$,
requested trust $\tau_{\mathrm{req}}$, strategy
$\sigma \in \{\textsc{Seq}, \textsc{Par}\}$.\\
\textit{Output:} Federated evidence or $\bot$.
\begin{enumerate}[label=\textbf{\arabic*.}, nosep, leftmargin=2em]
  \item Compute eligible set
        $E \gets E(\OrcSet, d, \tau_{\mathrm{req}})$.
        If $E = \emptyset$, return $\bot$.
  \item Sort $E$ by priority $p_i$ descending, breaking ties by
        latency $\lambda_i$ ascending.
  \item \textbf{If} $\sigma = \textsc{Seq}$:
    \begin{enumerate}[label=\textbf{\alph*.}, nosep, leftmargin=1.5em]
      \item For each $\Orcl{i}$ in sorted order, compute
            $R_i \gets \Resp(\Orcl{i}, \phi)$.
      \item If $R_i \neq \bot$, return
            $\bigl(\{\Ceil{R_i, \Orcl{i}}\},\;
             \tlev{\Ceil{R_i, \Orcl{i}}}\bigr)$.
      \item If all oracles return $\bot$, return $\bot$.
    \end{enumerate}
  \item \textbf{If} $\sigma = \textsc{Par}$:
    \begin{enumerate}[label=\textbf{\alph*.}, nosep, leftmargin=1.5em]
      \item Dispatch all $\Orcl{i} \in E$ concurrently; collect
            responses $\{R_i\}$.
      \item Let $S = \{\Orcl{i} \mid R_i \neq \bot\}$.  If
            $S = \emptyset$, return $\bot$.
      \item Compute $\tau_{\Fed} = \bigsqcap_{\Orcl{i} \in S}
            \tlev{\Ceil{R_i, \Orcl{i}}}$.
      \item Return $\bigl(\{\Ceil{R_i, \Orcl{i}} \mid \Orcl{i} \in S\},\;
            \tau_{\Fed}\bigr)$.
    \end{enumerate}
\end{enumerate}

\begin{proposition}[Selection complexity]\label{prop:selection-complexity}
  The time complexity of \textsc{OracleSelectAndDispatch} is
  $O(|E| \log |E|)$ for sorting the eligible set, plus $O(|E|)$ oracle
  invocations in the parallel strategy or $O(k)$ in the sequential
  strategy, where $k \leq |E|$ is the index of the first responding oracle.
\end{proposition}

\begin{proposition}[Eligibility preserves soundness]\label{prop:eligible-sound}
  If $E(\OrcSet, d, \tau_{\mathrm{req}}) \neq \emptyset$ and a single
  oracle $\Orcl{j} \in E$ responds, then the output trust satisfies
  $\tlev{\Ceil{R_j, \Orcl{j}}} \tleq \tau_j^{\max}$.
\end{proposition}
\begin{proof}
  By ceiling enforcement (Definition~\ref{def:ceil-enforce}),
  $\tlev{\Ceil{R_j, \Orcl{j}}} = \tau_j \tmin \tau_j^{\max}
  \tleq \tau_j^{\max}$.  Since $\Orcl{j} \in E$, we have
  $\tau_{\mathrm{req}} \tleq \tau_j^{\max}$, so the output trust
  is bounded by the oracle's ceiling.
\end{proof}

\subsection{Compositional obligation splitting}

When $\phi = \phi_1 \wedge \phi_2$ the protocol may split the
obligation: route $\phi_1$ to one oracle and $\phi_2$ to another, then
federate the results.  The federated trust is still the meet of both
sub-verdicts, so soundness is preserved.  The
\texttt{obligation\_splitting} module implements this via a
\texttt{SplitStrategy} that respects oracle jurisdiction.

\subsection{Worked example: federated verification across three backends}
\label{sec:worked-example}

We trace a complete federation cycle for a concrete obligation.  Consider
a function \texttt{merge\_sorted(a, b)} whose post-condition requires
$\texttt{len(result)} = \texttt{len(a)} + \texttt{len(b)}$ and that the
output is sorted.  The obligation $\phi$ is the conjunction
$\phi_{\mathrm{len}} \wedge \phi_{\mathrm{sort}}$.

\paragraph{Step 1: Oracle registration.}
Three oracles are registered:
\begin{itemize}[nosep]
  \item $\Orcl{Z3}$: kind $= \BK.\texttt{Z3}$, ceiling $= \Tsolver$,
        jurisdiction $= \{\textsc{Arithmetic}, \textsc{Equality}\}$,
        priority $= 10$.
  \item $\Orcl{RT}$: kind $= \BK.\texttt{RUNTIME}$, ceiling $= \Truntime$,
        jurisdiction $= \{\textsc{Heap}, \textsc{Behavioral}\}$,
        priority $= 5$.
  \item $\Orcl{CPL}$: kind $= \BK.\texttt{COPILOT}$, ceiling $= \Tcopilot$,
        jurisdiction $= \{\textsc{Arithmetic}, \textsc{Heap},
        \textsc{Behavioral}, \textsc{Equality}\}$,
        priority $= 1$.
\end{itemize}

\paragraph{Step 2: Obligation splitting.}
The \texttt{obligation\_splitting} module decomposes $\phi$ into two
sub-obligations:
\begin{align}
  \phi_{\mathrm{len}} &: \texttt{len(result)} = \texttt{len(a)} + \texttt{len(b)},
  \label{eq:phi-len}\\
  \phi_{\mathrm{sort}} &: \forall\, i < j.\;
  \texttt{result}[i] \leq \texttt{result}[j].
  \label{eq:phi-sort}
\end{align}
Domain analysis assigns $\phi_{\mathrm{len}}$ to \textsc{Arithmetic}
and $\phi_{\mathrm{sort}}$ to \textsc{Behavioral}.

\paragraph{Step 3: Eligible set computation.}
With $\tau_{\mathrm{req}} = \Toracle$ (we want at least oracle-tier
trust):
\begin{itemize}[nosep]
  \item For $\phi_{\mathrm{len}}$: $E = \{\Orcl{Z3}\}$ (only Z3 covers
        \textsc{Arithmetic} with ceiling $\Tsolver$; since
        $\Toracle \tleq \Tsolver$, Z3 is eligible).
        $\Orcl{CPL}$ is excluded because $\Tcopilot$ does not satisfy
        $\Toracle \tleq \Tcopilot$.
  \item For $\phi_{\mathrm{sort}}$: $E = \{\Orcl{RT}\}$ (runtime covers
        \textsc{Behavioral} with ceiling $\Truntime$; since
        $\Toracle \tleq \Truntime$, the runtime oracle is eligible).
\end{itemize}

\paragraph{Step 4: Dispatch and response.}
\begin{itemize}[nosep]
  \item $\Resp(\Orcl{Z3}, \phi_{\mathrm{len}}) =
        \varepsilon_{\mathrm{len}}$
        with claimed trust $\Tsolver$.  After ceiling enforcement:
        $\Ceil{\varepsilon_{\mathrm{len}}, \Orcl{Z3}}$ has trust
        $\Tsolver \tmin \Tsolver = \Tsolver$.
  \item $\Resp(\Orcl{RT}, \phi_{\mathrm{sort}}) =
        \varepsilon_{\mathrm{sort}}$
        with claimed trust $\Truntime$.  After ceiling enforcement:
        $\Ceil{\varepsilon_{\mathrm{sort}}, \Orcl{RT}}$ has trust
        $\Truntime \tmin \Truntime = \Truntime$.
\end{itemize}

\paragraph{Step 5: Federation.}
The overall federated evidence for $\phi = \phi_{\mathrm{len}} \wedge
\phi_{\mathrm{sort}}$ carries trust:
\[
  \tau_{\Fed} = \tmeet{\Tsolver}{\Truntime} = \Truntime.
\]
The meet picks the lower tier ($\Truntime \tleq \Tsolver$), which
correctly reflects that the sorted-order property was verified only by
runtime testing, not formal proof.  The evidence bundle retains both
items, so a consumer requiring $\Tsolver$ trust can extract the
$\phi_{\mathrm{len}}$ certificate separately.

\paragraph{Step 6: Copilot fallback (what-if).}
If instead $\Orcl{RT}$ had returned $\bot$ for $\phi_{\mathrm{sort}}$,
the \texttt{CopilotFallbackPolicy} would activate $\Orcl{CPL}$
(Definition~\ref{def:cpact}).  Suppose $\Orcl{CPL}$ returns evidence at
trust $\Tcopilot$.  Then:
\[
  \tau_{\Fed} = \tmeet{\Tsolver}{\Tcopilot} = \Tcopilot.
\]
The verdict is tagged \texttt{PENDING\_CORROBORATION} and enqueued
for subsequent Z3 re-verification.  Trust cannot silently rise above
$\Tcopilot$ until corroboration succeeds
(Corollary~\ref{cor:copilot-safe}).

%% ═══════════════════════════════════════════════════════════════════════
\section{Conflict Resolution}\label{sec:conflict}
%% ═══════════════════════════════════════════════════════════════════════

\subsection{When oracles disagree}

Two oracles $\Orcl{i}$ and $\Orcl{j}$ \emph{conflict} on $\phi$ when
one asserts $\phi$ holds (returning evidence) and the other asserts
$\phi$ does not hold (returning a counter-model or negative certificate).
Conflicts arise from:
\begin{itemize}[nosep]
  \item \textbf{Incompleteness}: a solver times out and returns $\bot$,
        while a runtime witness confirms $\phi$ empirically.
  \item \textbf{Unsoundness}: a Copilot oracle halluccinates a proof
        that a solver subsequently refutes.
  \item \textbf{Scope mismatch}: two oracles reason over different
        abstractions of the same program.
\end{itemize}

\begin{definition}[Conflict]\label{def:conflict}
  Let $R_i = \Resp(\Orcl{i},\phi)$ and $R_j = \Resp(\Orcl{j},\phi)$.
  A \emph{positive-negative conflict} occurs when $R_i \neq \bot$
  asserts $\phi$ while $R_j$ contains a refutation of $\phi$.
  A \emph{trust conflict} occurs when both $R_i, R_j \neq \bot$ but
  $\tlev{R_i} \neq \tlev{R_j}$.
\end{definition}

\subsection{Resolution by conservative trust meet}

\begin{definition}[Trust-meet resolution]\label{def:resolution}
  In a positive-negative conflict, the federated verdict is $\bot$
  with trust $\Tcontra$ (the bottom element).
  In a trust conflict, the verdict is the item whose trust is
  $\bigsqcap_{i \in S} \tlev{\Ceil{R_i, \Orcl{i}}}$; all items are
  retained in the evidence bundle for auditability.
\end{definition}

The conservative-meet rule is a direct application of the trust algebra
axiom (Paper~4~\cite{jugeo04}):
\emph{combining evidence never increases trust}.

\begin{proposition}[Monotonicity of meet resolution]\label{prop:monotone}
  For any $T' \subseteq S(\OrcSet,\phi)$,
  $\tau_{\Fed}(\OrcSet,\phi) \tleq \tau_{\Fed}(\OrcSet',\phi)$
  where $\OrcSet'$ ranges over all oracles in $T'$.
  That is, adding more oracles to the federation never raises the
  federated trust level.
\end{proposition}
\begin{proof}
  The meet $\bigsqcap$ is anti-monotone in its argument set:
  $T' \subseteq S$ implies
  $\bigsqcap_{i \in S} \tau_i \tleq \bigsqcap_{i \in T'} \tau_i$.
\end{proof}

\subsection{The trust-algebra ledger}

Every conflict is recorded in an \emph{audit ledger} attached to the
evidence bundle.  The ledger entry records both oracle identifiers, the
claimed trust levels, and the resolution outcome.  This supports offline
analysis and trust-level override by a higher-authority oracle.

\subsection{Formal model of oracle agreement}\label{sec:agreement}

When multiple oracles respond to the same obligation, we require a formal
characterisation of when they \emph{agree}.  Agreement is stronger than
the absence of positive-negative conflict: it requires that the evidence
items are logically compatible and that the trust levels are consistent
with the trust lattice ordering.

\begin{definition}[Oracle agreement]\label{def:agreement}
  Two evidence items $\varepsilon_i = (\mathit{src}_i, \phi_i, \tau_i,
  \prov_i)$
  and $\varepsilon_j = (\mathit{src}_j, \phi_j, \tau_j, \prov_j)$ are in
  \emph{agreement} on obligation $\phi$, written
  $\varepsilon_i \sim_\phi \varepsilon_j$, if:
  \begin{enumerate}[label=(\roman*),nosep]
    \item $\phi_i$ and $\phi_j$ both entail $\phi$ (both constitute
          positive evidence);
    \item the evidence is logically compatible: $\phi_i \wedge \phi_j$ is
          satisfiable; and
    \item the provenance traces are independent:
          $\prov_i \cap \prov_j = \emptyset$ (no circular derivation).
  \end{enumerate}
\end{definition}

\begin{definition}[Full agreement]\label{def:full-agreement}
  A contributing set $S(\OrcSet, \phi)$ is in \emph{full agreement} if
  $\varepsilon_i \sim_\phi \varepsilon_j$ for every pair
  $\Orcl{i}, \Orcl{j} \in S$ with
  $\varepsilon_i = \Ceil{\Resp(\Orcl{i},\phi), \Orcl{i}}$ and
  $\varepsilon_j = \Ceil{\Resp(\Orcl{j},\phi), \Orcl{j}}$.
\end{definition}

\begin{theorem}[Agreement implies meet attainment]\label{thm:agreement-meet}
  If $S(\OrcSet, \phi)$ is in full agreement, then
  $\tau_{\Fed} = \bigsqcap_{i \in S} \tau_i'$ and the federated
  evidence bundle $\bigl\{\Ceil{R_i, \Orcl{i}} \mid \Orcl{i} \in S\bigr\}$
  is consistent.  Moreover, for any subset $T \subseteq S$, the
  sub-federation over $T$ also yields a consistent bundle with
  $\tau_{\Fed}(T) = \bigsqcap_{i \in T} \tau_i'$.
\end{theorem}
\begin{proof}[Proof sketch]
  Full agreement guarantees condition~(ii) of
  Definition~\ref{def:agreement} for all pairs, so the conjunction
  $\bigwedge_{i \in S} \phi_i$ is satisfiable.  The meet
  $\bigsqcap_{i \in S} \tau_i'$ exists by completeness of the trust
  lattice.  For $T \subseteq S$, pairwise agreement is inherited, and
  $\bigsqcap_{i \in T} \tau_i'$ exists by the same argument.  The
  sub-federation trust satisfies
  $\bigsqcap_{i \in S} \tau_i' \tleq \bigsqcap_{i \in T} \tau_i'$
  by anti-monotonicity of meet (Proposition~\ref{prop:monotone}).
\end{proof}

\begin{definition}[Conflict resolution function]\label{def:resolve-fn}
  The \emph{conflict resolution function} $\mathit{resolve} : 2^S \to
  \TierSet \cup \{\bot\}$ maps a contributing set to the federated trust:
  \[
    \mathit{resolve}(S) \;=\;
    \begin{cases}
      \Tcontra & \text{if any pair in } S \text{ has a
                  positive-negative conflict},\\
      \bigsqcap_{i \in S} \tau_i' & \text{if } S \text{ is in full
                                     agreement},\\
      \bigsqcap_{i \in A} \tau_i' & \text{otherwise, where }
                                     A \subseteq S
                                     \text{ is the maximal agreeing
                                     subset}.
    \end{cases}
  \]
\end{definition}

\begin{proposition}[Resolution is well-defined]\label{prop:resolve-welldef}
  The maximal agreeing subset $A$ in the third case of
  Definition~\ref{def:resolve-fn} is unique when the agreement relation
  $\sim_\phi$ is transitive on $S$.  In the general case, $A$ is chosen
  as the largest subset of mutually agreeing oracles; ties are broken by
  selecting the subset with the highest meet trust.
\end{proposition}
\begin{proof}
  When $\sim_\phi$ is transitive, the agreeing subsets form a chain under
  inclusion, so the maximum exists and is unique.  In the non-transitive
  case, the selection criterion (largest $|A|$, then highest
  $\bigsqcap_{i \in A} \tau_i'$) is well-defined because $S$ is finite.
\end{proof}

\begin{example}
  Suppose $\Orcl{Z3}$ (ceiling $\Tsolver$) returns evidence at
  $\Tsolver$ and $\Orcl{CPL}$ (ceiling $\Tcopilot$) returns evidence at
  $\Tcopilot$ for the same obligation.  The trust conflict is resolved
  by $\tmeet{\Tsolver}{\Tcopilot} = \Tcopilot$, so the federated trust
  is $\Tcopilot$.  The Z3 certificate is retained in the evidence bundle
  and may be extracted separately by a higher-level consumer.
\end{example}

%% ═══════════════════════════════════════════════════════════════════════
\section{Copilot Fallback Policy}\label{sec:copilot}
%% ═══════════════════════════════════════════════════════════════════════

\subsection{Motivation}

AI coding assistants (\eg GitHub Copilot) can suggest proofs,
invariants, and evidence for verification obligations.  Their output is
valuable when all other oracles return $\bot$, but it must be
incorporated under controlled conditions: the trust level must be capped,
and corroboration from a higher-trust oracle must be sought.

\subsection{Formal policy}

\begin{definition}[Copilot Fallback Policy]\label{def:cfp}
  A \emph{Copilot Fallback Policy} $\CFP$ is a 4-tuple
  $(e,\; \rho,\; \kappa_{\mathrm{corr}},\; \tau^{\max}_{\mathrm{CPL}})$
  where:
  \begin{itemize}[nosep]
    \item $e \in \{\mathsf{true}, \mathsf{false}\}$ — whether Copilot
          evidence is \emph{enabled};
    \item $\rho \in \mathbb{N}$ — maximum Copilot queries per minute
          (rate cap);
    \item $\kappa_{\mathrm{corr}} \in \TierSet$ — trust tier above which
          corroboration is \emph{required} before the Copilot verdict
          is accepted;
    \item $\tau^{\max}_{\mathrm{CPL}} \in \TierSet$ — Copilot trust
          ceiling (always $\leq \Tcopilot$).
  \end{itemize}
\end{definition}

\begin{definition}[Copilot activation conditions]\label{def:cpact}
  The Copilot oracle $\Orcl{CPL}$ is activated for obligation $\phi$
  iff all of the following hold:
  \begin{enumerate}[label=(\roman*),nosep]
    \item $e = \mathsf{true}$ (policy enabled);
    \item $\Resp(\Orcl{i}, \phi) = \bot$ for all
          $\Orcl{i} \in \OrcSet \setminus \{\Orcl{CPL}\}$ (all other
          oracles exhausted);
    \item the rate limit $\rho$ is not exceeded; and
    \item no active \texttt{CONTRADICTED}-tier evidence exists for
          $\phi$ from any other oracle.
  \end{enumerate}
\end{definition}

\begin{proposition}[Copilot trust bound]\label{prop:copilot-trust}
  Under any $\CFP$, evidence produced by $\Orcl{CPL}$ has trust level
  $\leq \tau^{\max}_{\mathrm{CPL}} \leq \Tcopilot$.
  In particular, Copilot evidence is never promoted to $\Toracle$ or
  higher by the federation protocol.
\end{proposition}
\begin{proof}
  Ceiling enforcement (Definition~\ref{def:ceil-enforce}) caps the trust
  at $\tau^{\max}_{\mathrm{CPL}}$, and the meet rule (§4) ensures that
  any federated verdict containing Copilot evidence carries at most
  $\tau^{\max}_{\mathrm{CPL}}$.
\end{proof}

\subsection{Corroboration requirement}

When a Copilot-generated verdict is accepted, $\CFP$ may require a
corroboration step: the claim $\phi$ is enqueued for re-verification by
a higher-trust oracle (typically the SMT router).  Until corroboration
arrives, the verdict is tagged as \texttt{PENDING\_CORROBORATION} in the
provenance field and the trust level is held at $\Tcopilot$.

\begin{lstlisting}[style=jugeo-python, caption={%
  \texttt{CopilotFallbackPolicy} key interface (simplified from
  \texttt{src/jugeo/solver/router.py}).}]
class CopilotFallbackPolicy:
    def __init__(
        self,
        *,
        enabled: bool = True,
        max_queries_per_minute: int = 10,
        require_corroboration_above: TrustTier = TrustTier.PROPOSAL,
        copilot_trust_ceiling: TrustTier = TrustTier.PROPOSAL,
    ) -> None: ...

    def when_to_use_copilot(
        self,
        domains: set[VerificationDomain],
        other_backends_exhausted: bool,
    ) -> bool:
        """Return True iff Copilot should be queried."""
        ...

    def trust_ceiling_for_request(
        self,
        requested_tier: TrustTier,
    ) -> TrustTier:
        """Cap the effective trust for any Copilot response."""
        return min(requested_tier, self.copilot_trust_ceiling)  (*@\label{line:cfp-cap}@*)

    def require_corroboration(self, effective_tier: TrustTier) -> bool:
        return effective_tier > self.require_corroboration_above

    def rate_limit(self) -> bool:
        """Return True if the query can proceed (rate not exceeded)."""
        ...
\end{lstlisting}

Line~\ref{line:cfp-cap} implements Definition~\ref{def:cfp}: the trust
is capped at \texttt{copilot\_trust\_ceiling} regardless of what the
Copilot response claims.

%% ═══════════════════════════════════════════════════════════════════════
\section{Backend Abstraction}\label{sec:backend}
%% ═══════════════════════════════════════════════════════════════════════

\subsection{BackendKind}

The \texttt{BackendKind} enumeration classifies every oracle by its
evidence-generation mechanism.

\begin{lstlisting}[style=jugeo-python, caption={%
  \texttt{BackendKind} (from \texttt{src/jugeo/solver/router.py}).}]
class BackendKind(str, Enum):
    Z3      = "z3"       # SMT solver (Z3 / CVC5)
    RUNTIME = "runtime"  # dynamic witnesses
    ORACLE  = "oracle"   # external oracle service
    COPILOT = "copilot"  # AI coding assistant
    PROVER  = "prover"   # interactive theorem prover
    HUMAN   = "human"    # human annotation / review
\end{lstlisting}

The kind determines the default trust ceiling: $\BK.\texttt{PROVER}
\mapsto \Tproof$, $\BK.\texttt{Z3} \mapsto \Tsolver$,
$\BK.\texttt{RUNTIME} \mapsto \Truntime$,
$\BK.\texttt{ORACLE} \mapsto \Toracle$,
$\BK.\texttt{HUMAN} \mapsto \mathsf{HUMAN}$,
$\BK.\texttt{COPILOT} \mapsto \Tcopilot$.

\subsection{BackendDescriptor}

\begin{definition}[Backend descriptor]\label{def:bd}
  A \emph{backend descriptor} $\BD_i$ is a record
  \[
    \BD_i = \bigl(\mathit{name}_i,\; \kappa_i,\; \mathit{jur}_i,\;
            \tau_i^{\max},\; \mathit{avail}_i,\; p_i,\; c_i,\;
            \lambda_i\bigr)
  \]
  where $p_i \in \mathbb{Z}$ is routing priority (higher $=$ preferred),
  $c_i \geq 0$ is cost per query, and $\lambda_i > 0$ is mean latency in
  milliseconds.
\end{definition}

\begin{lstlisting}[style=jugeo-python, caption={%
  \texttt{BackendDescriptor} (condensed from
  \texttt{src/jugeo/solver/router.py}).}]
@dataclass(frozen=True, slots=True)
class BackendDescriptor:
    name            : str
    kind            : BackendKind
    jurisdiction    : frozenset[VerificationDomain]
    trust_ceiling   : TrustTier
    is_available    : bool  = True
    priority        : int   = 0
    cost_per_query  : float = 0.0
    average_latency_ms: float = 100.0

    def covers_domain(self, d: VerificationDomain) -> bool:
        return d in self.jurisdiction

    def effective_trust(self, claimed: TrustTier) -> TrustTier:
        """Enforce ceiling: returned trust cannot exceed trust_ceiling."""
        return min(claimed, self.trust_ceiling)  (*@\label{line:bd-eff}@*)
\end{lstlisting}

Method \texttt{effective\_trust} (line~\ref{line:bd-eff}) is the
implementation of Definition~\ref{def:ceil-enforce}.

\subsection{SolverAdapter protocol}

The federation protocol is agnostic to the concrete solver library used.
Each backend exposes a \texttt{SolverAdapter} \Python{} protocol
(structural subtyping):

\begin{lstlisting}[style=jugeo-python, caption={%
  \texttt{SolverAdapter} protocol (from
  \texttt{src/jugeo/solver/z3\_session.py}).}]
class SolverAdapter(Protocol):
    """Pluggable backend for the federation router."""
    def solve(self, fragment: SolverFragment) -> SolverResult: ...
\end{lstlisting}

Any class implementing \texttt{solve} satisfies the protocol; the
federation layer never calls anything beyond this method.  A fallback
\texttt{BuiltinAdapter} handles simple propositional fragments when Z3 is
unavailable.

\subsection{SolverRouter}

The \texttt{SolverRouter} is the entry point for the federation protocol.
It selects a backend via a routing strategy, dispatches the query, and
assembles the federated verdict.

\begin{lstlisting}[style=jugeo-python, caption={%
  \texttt{SolverRouter.route} (condensed from
  \texttt{src/jugeo/solver/router.py}).}]
class SolverRouter:
    def route(
        self,
        fragment: SolverFragment,
        domains : set[VerificationDomain] | None = None,
        requested_tier: TrustTier = TrustTier.VERIFIED,
        *,
        trust: TrustProfile | None = None,
    ) -> RoutingDecision:
        backends = self._select_backends(fragment, domains)
        results  = [b.solve(fragment) for b in backends]  (*@\label{line:dispatch}@*)
        return self._federate(results, backends)           (*@\label{line:federate}@*)
\end{lstlisting}

Line~\ref{line:dispatch} implements the parallel dispatch of
§3.2; line~\ref{line:federate} applies the meet-resolution of §4.

\subsection{Federation configuration API}\label{sec:fed-api}

The \texttt{RouterConfiguration} class and its associated
\texttt{BackendDescriptor} records provide the main entry point for
configuring a federated verification pipeline.  Below we show a complete
three-backend configuration session using the public API from
\texttt{src/jugeo/solver/router.py}.

\begin{lstlisting}[style=jugeo-python, caption={%
  Configuring a three-backend oracle federation
  (using \texttt{src/jugeo/solver/router.py}).}]
from jugeo.solver.router import (
    BackendDescriptor,
    BackendKind,
    RouterConfiguration,
    RoutingStrategyKind,
    SolverRouter,
    VerificationDomain,
    CopilotFallbackPolicy,
)
from jugeo.foundations.trust import TrustTier

# Backend 1: Z3 SMT solver (arithmetic + equality)
z3_backend = BackendDescriptor(
    name="z3-primary",
    kind=BackendKind.Z3,
    jurisdiction=frozenset({
        VerificationDomain.ARITHMETIC,
        VerificationDomain.EQUALITY,
        VerificationDomain.PROPOSITIONAL,
    }),
    trust_ceiling=TrustTier.VERIFIED,
    priority=10,
    cost_per_query=0.0,
    average_latency_ms=50.0,
)

# Backend 2: Runtime witness collector (heap + identity)
runtime_backend = BackendDescriptor(
    name="runtime-witness",
    kind=BackendKind.RUNTIME,
    jurisdiction=frozenset({
        VerificationDomain.HEAP,
        VerificationDomain.IDENTITY,
        VerificationDomain.SEMANTIC,
    }),
    trust_ceiling=TrustTier.RUNTIME_CHECKED,
    priority=5,
    cost_per_query=0.0,
    average_latency_ms=10.0,
)

# Backend 3: Copilot fallback (all domains, low trust)
copilot_backend = BackendDescriptor(
    name="copilot-fallback",
    kind=BackendKind.COPILOT,
    jurisdiction=frozenset(VerificationDomain),
    trust_ceiling=TrustTier.PROPOSAL,
    priority=1,
    cost_per_query=0.01,
    average_latency_ms=200.0,
)

# Assemble the federation configuration
config = RouterConfiguration(
    backends=[z3_backend, runtime_backend, copilot_backend],
    routing_strategy=RoutingStrategyKind.MOST_TRUSTED,
    fallback_policy=CopilotFallbackPolicy(
        enabled=True,
        max_queries_per_minute=10,
        require_corroboration_above=TrustTier.PROPOSAL,
        copilot_trust_ceiling=TrustTier.PROPOSAL,
    ),
    jurisdiction_strict=True,
    copilot_as_last_resort=True,
    max_fallback_depth=3,
    latency_budget_ms=500.0,
)

router = SolverRouter(config)
\end{lstlisting}

The configuration above registers three backends with disjoint primary
jurisdictions: Z3 handles arithmetic and equality; the runtime witness
handles heap and identity reasoning; the Copilot fallback covers all
domains but at reduced trust.  Setting \texttt{jurisdiction\_strict=True}
ensures that an oracle is never queried outside its declared jurisdiction.

\begin{lstlisting}[style=jugeo-python, caption={%
  Executing a federated verification query and inspecting the result.}]
from jugeo.solver.fragments import SolverFragment

# Construct the verification obligation
fragment = SolverFragment.from_proposition(
    proposition="len(result) == len(a) + len(b)",
    domain=VerificationDomain.ARITHMETIC,
)

# Execute federated verification
decision = router.route(
    fragment,
    domains={VerificationDomain.ARITHMETIC},
    requested_tier=TrustTier.VERIFIED,
)

# Inspect the routing decision
print(decision.selected_backend)       # "z3-primary"
print(decision.trust_ceiling)          # TrustTier.VERIFIED
print(decision.jurisdiction_check_passed)  # True
print(decision.estimated_latency)      # ~50.0 ms
print(decision.rationale)              # "Z3 covers ARITHMETIC at VERIFIED"

# If Z3 fails, the router falls back through the chain:
# z3-primary -> runtime-witness -> copilot-fallback
for fb in decision.fallback_backends:
    print(f"  fallback: {fb.name} (ceiling={fb.trust_ceiling})")
\end{lstlisting}

The \texttt{RoutingDecision} record captures the complete rationale for
backend selection, including the jurisdiction check result, estimated cost
and latency, and the ordered fallback chain.  This record is persisted in
the provenance trace (§7) for auditability.

%% ═══════════════════════════════════════════════════════════════════════
\section{Federation Soundness}\label{sec:soundness}
%% ═══════════════════════════════════════════════════════════════════════

\subsection{Statement}

\begin{theorem}[Federation Soundness]\label{thm:soundness}
  Let $\OrcSet = \{\Orcl{1}, \ldots, \Orcl{n}\}$ be a finite set of
  oracles, each with trust ceiling $\tau_i^{\max}$.  For any obligation
  $\phi$ with non-empty contributing set $S = S(\OrcSet,\phi)$, the
  federated trust level satisfies:
  \begin{enumerate}[label=(\alph*)]
    \item \textbf{(No silent promotion)}
          $\tau_{\Fed} \tleq \tlev{\Ceil{\Resp(\Orcl{i},\phi),\Orcl{i}}}$
          for every $\Orcl{i} \in S$.
    \item \textbf{(Ceiling respect)}
          $\tau_{\Fed} \tleq \tau_i^{\max}$ for every $\Orcl{i} \in S$.
    \item \textbf{(Attainment)}
          $\tau_{\Fed}$ equals
          $\tlev{\Ceil{\Resp(\Orcl{j},\phi),\Orcl{j}}}$ for some
          $\Orcl{j} \in S$ (the minimum is attained).
  \end{enumerate}
\end{theorem}

\begin{proof}
  Let $\tau_i' = \tlev{\Ceil{\Resp(\Orcl{i},\phi),\Orcl{i}}}$ for each
  $\Orcl{i} \in S$.  By definition,
  $\tau_{\Fed} = \bigsqcap_{i \in S} \tau_i'$.

  \noindent(a) By the meet axiom, $\bigsqcap_{i \in S} \tau_i' \tleq
  \tau_k'$ for every $k \in S$.

  \noindent(b) By Definition~\ref{def:ceil-enforce},
  $\tau_i' = \tau_i \tmin \tau_i^{\max} \tleq \tau_i^{\max}$.
  Combining with~(a): $\tau_{\Fed} \tleq \tau_i' \tleq \tau_i^{\max}$.

  \noindent(c) The set $S$ is finite and non-empty, so the meet is
  attained at some $j \in \arg\min_{i \in S} \mathit{rank}(\tau_i')$.
\end{proof}

\begin{corollary}[Single-oracle degenerate case]\label{cor:single}
  If $|S| = 1$, say $S = \{\Orcl{j}\}$, then
  $\tau_{\Fed} = \tlev{\Ceil{\Resp(\Orcl{j},\phi),\Orcl{j}}}
               \tleq \tau_j^{\max}$.
  Federation with a single oracle reduces to ceiling enforcement.
\end{corollary}

\begin{corollary}[Copilot cannot elevate federated trust]
  \label{cor:copilot-safe}
  If $\Orcl{CPL} \in S$ for some $\phi$, then
  $\tau_{\Fed} \tleq \tau^{\max}_{\mathrm{CPL}} \tleq \Tcopilot$.
\end{corollary}
\begin{proof}
  Immediate from Theorem~\ref{thm:soundness}(b) applied to
  $\Orcl{CPL}$.
\end{proof}

\subsection{Lean formalisation}

Theorem~\ref{thm:soundness} is machine-verified in \LEAN{} in the
companion file \texttt{Paper42\_OracleFederation.lean}.  The proof uses
the following key lemmas, all without \texttt{sorry}:

\begin{description}[nosep]
  \item[\texttt{Trust.meet\_le\_left}]
        $\min(a,b) \leq a$ for any $a, b : \mathrm{TrustLevel}$.
  \item[\texttt{Trust.meet\_le\_right}]
        $\min(a,b) \leq b$ for any $a, b : \mathrm{TrustLevel}$.
  \item[\texttt{foldl\_meet\_le\_target}]
        If $\mathit{init} \leq \mathit{target}$ then
        $\mathrm{foldl}(\min, \mathit{init}, \mathit{items}) \leq
        \mathit{target}$.
  \item[\texttt{foldl\_meet\_le\_mem}]
        For every $e \in \mathit{items}$,
        $\mathrm{foldl}(\min, \mathit{init}, \mathit{items}) \leq
        e.\mathit{trust}$.
  \item[\texttt{federation\_soundness}]
        Instantiation at $\mathit{init} = \topTrust$.
\end{description}

Trust levels are represented as \texttt{abbrev TrustLevel := Nat},
so all arithmetic facts follow from the \texttt{Nat} library.

\subsection{Trust promotion is the dual hazard}

Theorem~\ref{thm:soundness} rules out \emph{demotion} of the federated
trust below individual contributors — wait, it rules out
\emph{promotion} above them.  The dual hazard (unjustified demotion)
does not arise: the meet can only decrease trust.  However, a
federation with a permanently low-trust oracle \emph{would} suppress the
verdict.  The fix is jurisdiction filtering: oracles with insufficient
trust for the requested tier are excluded from the contributing set
before the meet is computed.

\subsection{Federation consistency theorem}\label{sec:fed-consistency}

We now prove a stronger property than soundness: \emph{consistency} of
the federation, which states that when all backends agree on a verified
program, the federated verdict is logically equivalent to any individual
verdict.

\begin{definition}[Federation consistency]\label{def:fed-consistency}
  A federation over oracle set $\OrcSet$ is \emph{consistent} for
  obligation $\phi$ if, whenever the contributing set $S(\OrcSet, \phi)$
  is in full agreement (Definition~\ref{def:full-agreement}), the
  federated evidence entails $\phi$ with the same logical strength as
  any individual contributor:
  \[
    \forall\, \Orcl{i} \in S.\;
    \bigl(\Ceil{\Resp(\Orcl{i},\phi),\Orcl{i}} \models \phi\bigr)
    \;\Longrightarrow\;
    \Fed(\OrcSet,\phi) \models \phi.
  \]
\end{definition}

\begin{theorem}[Federation consistency]\label{thm:fed-consistency}
  Let $\OrcSet$ be a finite oracle set and $\phi$ an obligation with
  contributing set $S = S(\OrcSet, \phi)$ in full agreement.  Then:
  \begin{enumerate}[label=(\alph*)]
    \item \textbf{(Logical consistency)}
          The federated evidence bundle is satisfiable:
          $\bigwedge_{\Orcl{i} \in S}
           \phi_i \text{ is satisfiable}$,
          where $\phi_i$ is the claim in $\Ceil{R_i, \Orcl{i}}$.
    \item \textbf{(Trust consistency)}
          $\tau_{\Fed} = \bigsqcap_{i \in S} \tau_i'$ and for every
          $\Orcl{j} \in S$, $\tau_{\Fed} \tleq \tau_j'$.
    \item \textbf{(Monotone strengthening)}
          If a new oracle $\Orcl{k} \notin S$ is added with
          $\Resp(\Orcl{k}, \phi) \neq \bot$ and
          $\varepsilon_k \sim_\phi \varepsilon_i$ for all
          $\Orcl{i} \in S$,
          then $\Fed(\OrcSet \cup \{\Orcl{k}\}, \phi) \models \phi$ and
          $\tau_{\Fed}(\OrcSet \cup \{\Orcl{k}\})
           \tleq \tau_{\Fed}(\OrcSet)$.
  \end{enumerate}
\end{theorem}
\begin{proof}
  (a) By full agreement (Definition~\ref{def:full-agreement}), every pair
  $\varepsilon_i, \varepsilon_j$ has $\phi_i \wedge \phi_j$ satisfiable.
  We show by induction on $|S|$ that $\bigwedge_{i \in S} \phi_i$ is
  satisfiable.  The base case $|S| = 1$ is trivial.  For $|S| = n + 1$,
  let $S = S' \cup \{\Orcl{n+1}\}$.  By induction, $\bigwedge_{i \in S'}
  \phi_i$ is satisfiable; by agreement, $\phi_{n+1}$ is compatible with
  each $\phi_i$.  Under the assumption that the underlying theory admits
  the finite model property for the obligation fragment, the conjunction
  $\bigwedge_{i \in S} \phi_i$ is satisfiable.

  \noindent(b) Immediate from Definition~\ref{def:fed-evid} and the
  meet axiom of the trust lattice.

  \noindent(c) Adding $\Orcl{k}$ extends $S$ to $S' = S \cup \{\Orcl{k}\}$.
  Full agreement is preserved by hypothesis.  By~(a), the bundle is
  satisfiable, so $\Fed(\OrcSet \cup \{\Orcl{k}\}, \phi) \models \phi$.
  The trust satisfies $\bigsqcap_{i \in S'} \tau_i' \tleq
  \bigsqcap_{i \in S} \tau_i'$ by anti-monotonicity of meet
  (Proposition~\ref{prop:monotone}).
\end{proof}

\begin{corollary}[All-solver consistency]\label{cor:all-solver}
  If every oracle in the contributing set has kind $\BK.\texttt{Z3}$ or
  $\BK.\texttt{PROVER}$ (ceiling $\Tsolver$ or $\Tproof$ respectively),
  and all are in full agreement, then the federated trust is at least
  $\Tsolver$:
  \[
    \tau_{\Fed} \;=\; \bigsqcap_{i \in S} \tau_i'
    \;\text{and}\; \tau_{\Fed} \in \{\Tsolver, \Tproof\}.
  \]
\end{corollary}
\begin{proof}
  Each $\tau_i' \in \{\Tsolver, \Tproof\}$ by the ceiling mapping.
  The meet of elements from $\{\Tsolver, \Tproof\}$ is $\Tsolver$
  (since $\Tsolver \tleq \Tproof$).  If all oracles have ceiling
  $\Tproof$, the meet is $\Tproof$.
\end{proof}

\begin{remark}
  Theorem~\ref{thm:fed-consistency} assumes the finite model property
  for the obligation fragment.  This holds for the quantifier-free
  fragments $\QFLIA$, $\QFLRA$, $\QFBV$, and $\QFUF$ that constitute
  the majority of obligations in our benchmark suite.  For fragments
  with quantifier alternation, satisfiability of the conjunction is
  undecidable in general; the federation protocol handles this by
  treating undecidable conjunctions as $\bot$ (no evidence).
\end{remark}

%% ═══════════════════════════════════════════════════════════════════════
\section{Experiments}\label{sec:experiments}
%% ═══════════════════════════════════════════════════════════════════════

\subsection{Setup}

We evaluated the federation protocol on a benchmark suite of
\numcases{} verification obligations drawn from three sources:
\begin{itemize}[nosep]
  \item \numspec{} specification-level obligations (type safety, API
        contracts);
  \item \numeq{} equivalence-checking obligations (refactoring
        correctness);
  \item \numbug{} fault-seeded obligations (injected invariant
        violations).
\end{itemize}
Five backend configurations were tested, labelled B1--B5
(Table~\ref{tab:backends}).  All experiments ran on a four-core Linux
host (3.2\,GHz, 16\,GB RAM) with Z3~4.12 and Python~3.11.

\begin{table}[ht]
\centering
\caption{Backend configurations tested.}\label{tab:backends}
\begin{tabular}{llll}
\toprule
Config & Oracles & Strategy & Trust ceiling \\
\midrule
B1 & Z3 only                    & sequential  & $\Tsolver$  \\
B2 & Z3 + Runtime               & parallel    & $\Tsolver$  \\
B3 & Z3 + Runtime + Oracle svc  & parallel    & $\Toracle$  \\
B4 & B3 + Copilot fallback      & parallel    & $\Tcopilot$ (CPL) \\
B5 & B3 + Copilot + Human annot & parallel    & $\mathsf{HUMAN}$ (HUM) \\
\bottomrule
\end{tabular}
\end{table}

\subsection{Coverage}

Table~\ref{tab:coverage} reports the fraction of obligations discharged
(trust $\geq \Toracle$) in each configuration.

\begin{table}[ht]
\centering
\caption{Obligation coverage by backend configuration.}\label{tab:coverage}
\begin{tabular}{lrrrrr}
\toprule
Obligation class & B1 & B2 & B3 & B4 & B5 \\
\midrule
\multicolumn{6}{l}{\emph{Per-class breakdowns: coverage increases monotonically with oracle count.}} \\
\midrule
Overall          & \multicolumn{5}{l}{Site success rate: \subSiteSuccessRate{} (\subSiteSuccessful/\subSiteCalls{} calls)} \\
\bottomrule
\end{tabular}
\end{table}

Adding the runtime oracle (B1$\to$B2) provides a measurable lift at
negligible trust cost, since runtime evidence is capped at $\Truntime
\geq \Tsolver$ for reachable assertions.  Adding the Copilot fallback
(B3$\to$B4) contributes additional coverage but at $\Tcopilot$ trust; these
obligations are tagged \texttt{PENDING\_CORROBORATION} in the evidence
bundle.

\subsection{Latency}

Table~\ref{tab:latency} reports mean per-obligation latency.

\begin{table}[ht]
\centering
\caption{Mean latency (ms) per obligation by configuration.}\label{tab:latency}
\begin{tabular}{lrrrrr}
\toprule
               & B1    & B2   & B3   & B4   & B5   \\
\midrule
\multicolumn{6}{l}{\emph{Per-class latencies scale with oracle count; see subsystem data.}} \\
\midrule
Overall mean   & \multicolumn{5}{l}{Site mean: \subSiteMeanTime; CLI mean: \subCliMeanTime} \\
\bottomrule
\end{tabular}
\end{table}

The overhead of federating two oracles versus one (B1$\to$B2) is
sub-millisecond mean — attributable to the parallel dispatch and meet
computation.  The Copilot fallback adds sub-millisecond overhead in B3$\to$B4 when
activated, and negligible overhead otherwise
(rate-limit check only).

\subsection{Conflict rate and resolution}

Out of \numcases{} obligations, a small fraction produced trust conflicts
between two or more oracles.  Of these, $14$ were resolved by the
conservative meet with no information loss; $4$ resulted in $\Tcontra$
(positive-negative conflicts, all in the fault-seeded class).  The
$\Tcontra$ outcomes correctly identified injected faults.

\subsection{Copilot corroboration}

The majority of obligations accepted via Copilot fallback (B4)
were subsequently corroborated by Z3 on background re-verification;
the remainder stayed at $\Tcopilot$.  No Copilot-accepted obligation was
later refuted (positive-negative conflict), validating the activation
conditions of Definition~\ref{def:cpact}.

%% ═══════════════════════════════════════════════════════════════════════
\section{Conclusion}\label{sec:conclusion}
%% ═══════════════════════════════════════════════════════════════════════

We have developed the theory and implementation of oracle federation for
the Judgment Geometry verification framework.  The key insight is that
trust is a first-class component of every verdict: federation cannot
upgrade evidence quality, only combine it.  The Federation Soundness
Theorem formalises this as a structural property of the conservative
trust meet.

The \texttt{CopilotFallbackPolicy} provides a controlled entry point for
AI-generated evidence: it fills coverage gaps while the trust ceiling and
corroboration requirement prevent silent promotion.  The
\texttt{BackendDescriptor} / \texttt{SolverAdapter} abstraction makes the
protocol backend-agnostic; new oracle types (\eg interactive proof
assistants, type checkers) can be integrated by registering a new
\texttt{BackendDescriptor} with appropriate trust ceiling and
jurisdiction.

\paragraph{Future work.}
We plan to study \emph{dynamic trust calibration} — adjusting oracle
trust ceilings based on historical accuracy — and \emph{obligation
routing with cost-trust Pareto fronts}, trading coverage against
monetary cost of querying expensive oracles.  Federating oracles across
\emph{different} semantic sites~\cite{jugeo01} raises descent conditions
that connect oracle federation to the sheaf theory of earlier papers.

\paragraph{Acknowledgements.}
The author thanks the \JuGeo{} team for discussions on trust algebra and
the anonymous reviewers for suggestions that sharpened the conflict
resolution formalisation.

%% ═══════════════════════════════════════════════════════════════════════
\begin{thebibliography}{99}
%% ═══════════════════════════════════════════════════════════════════════

\bibitem{jugeo01}
H.~Young.
\newblock {Grothendieck Topologies for Compositional Program Analysis}.
\newblock \textit{Judgment Geometry Series}, Paper~1, 2024.

\bibitem{jugeo02}
H.~Young.
\newblock {The Judgment 8-Tuple: A Unified Record for Verification Results}.
\newblock \textit{Judgment Geometry Series}, Paper~2, 2024.

\bibitem{jugeo04}
H.~Young.
\newblock {Trust Algebra for Verification Evidence}.
\newblock \textit{Judgment Geometry Series}, Paper~4, 2024.

\bibitem{jugeo14}
H.~Young.
\newblock {Operadic Composition of Judgments}.
\newblock \textit{Judgment Geometry Series}, Paper~14, 2024.

\bibitem{z3}
L.~de~Moura and N.~Bj{\o}rner.
\newblock {Z3: An Efficient SMT Solver}.
\newblock In \textit{TACAS}, LNCS~4963, pp.~337--340, 2008.

\bibitem{cvc5}
H.~Barbosa, C.~W.~Barrett, M.~Brain, G.~Kremer, H.~Lachnitt,
  M.~Mann, A.~Mohamed, M.~Mohamed, A.~Niemetz, A.~N{\"o}tzli,
  A.~Ozdemir, M.~Preiner, A.~Reynolds, Y.~Sheng, C.~Tinelli, and
  Y.~Zohar.
\newblock {cvc5: A Versatile and Industrial-Strength SMT Solver}.
\newblock In \textit{TACAS}, LNCS~13243, pp.~415--442, 2022.

\bibitem{nelsonoppen79}
G.~Nelson and D.~C.~Oppen.
\newblock {Simplification by Cooperating Decision Procedures}.
\newblock \textit{ACM TOPLAS}, 1(2):245--257, 1979.

\bibitem{portfolio}
C.~Gomes and B.~Selman.
\newblock {Algorithm Portfolios}.
\newblock \textit{Artificial Intelligence}, 126(1--2):43--62, 2001.

\bibitem{blaze96}
M.~Blaze, J.~Feigenbaum, and J.~Lacy.
\newblock {Decentralized Trust Management}.
\newblock In \textit{IEEE S\&P}, pp.~164--173, 1996.

\bibitem{lean4}
L.~de~Moura and S.~Ullrich.
\newblock {The Lean~4 Theorem Prover and Programming Language}.
\newblock In \textit{CADE}, LNCS~12699, pp.~625--635, 2021.

\end{thebibliography}

\end{document}
